% from main.tex
%%%%%%%%%%  scan idx 270  |  folio 257  %%%%%%%%%%
extérieur à \([a,b]\).
\item \(\parfamp_{f_i}(E_i)\subset[a_i,b_i]\) pour \(i=1,2\)
\[
 \Longrightarrow\parfamp_f(E_1\otimes_S^{\mathbf L}E_2)\subset[a,b].
\]
\end{enumerate}

\emph{Preuve.} Conjuguer (1.9) et (3.5).

\textbf{Corollaire 4.7.1.} Sous les hypothèses de (4.6), si, pour \(i=1,2\), \(f_i\) est de dimension parfaite \(\leq n_i\), \(f\) est de dimension parfaite \(\leq n_1+n_2\). Par suite, si \(f_1\) et \(f_2\) sont parfaits, il en est de même de \(f\).

\textbf{Corollaire 4.7.2.} Soit
\[
\begin{tikzcd}
X \arrow[d,"f"'] & X' \arrow[l,"g"'] \arrow[d,"f'"]\\
Y & Y' \arrow[l,"h"']
\end{tikzcd}
\]
un carré cartésien de schémas, où \(f\) est localement de type fini et \(X\) et \(Y'\) sont supposés tor-indépendants sur \(Y\), et soient \(E\in\Ob D(X)\), et \([a,b]\) un intervalle de \(\mathbb Z\). On a :
\[
 \parfamp_f(E)\subset[a,b]\Longrightarrow\parfamp_{f'}(\mathbf Lg^*(E))\subset[a,b],
\]
et par suite le foncteur \(\mathbf Lg^*:D^-(X)\to D^-(X')\) envoie \(D(f)_{\parf}\) (resp. \(D(f)_{\uparf}\)) dans \(D(f')_{\parf}\) (resp. \(D(f')_{\uparf}\)) (4.2).

\textbf{Proposition 4.8.} Soient \(S\) un schéma, \(f:X\to Y\) un \(S\)-morphisme propre de \(S\)-schémas localement de type fini, \(E\in\Ob D^b(X)\). Alors, sous réserve de validité de (2.1) pour \(f\) (donc en tout cas si \(S\) est localement noethérien ou si \(f\) est projectif) :
\begin{enumerate}[label=\alph*)]
\item Si \(E\) est parfait rel. à \(S\), il en est de même de \(Rf_*(E)\).
\item Supposons \(S\) quasi-compact, \(X\) et \(Y\) de type fini sur \(S\), soient \([a,b]\) un intervalle de \(\mathbb Z\), et \(N\) un entier tel que \(R^if_*(M)=0\) pour tout \(i>N\) et
